\documentclass{article} %210 mm × 297 mm
\usepackage[utf8]{inputenc}
\usepackage[a4paper,left=1.5cm, right=1.5cm, top=2cm, bottom=2cm]{geometry}
\usepackage{graphicx}
%\usepackage{blindtext}
\usepackage{multirow}
\usepackage{mhchem}
\usepackage{url}
\usepackage{subfigure}
\usepackage{amsfonts,latexsym} % para tener disponibilidad de diversos simbolos
\usepackage{enumerate}
%\usepackage{booktabs}
\usepackage{float}
%\usepackage{threeparttable}
\usepackage{array,colortbl}
%\usepackage{ifpdf}
%\usepackage{rotating}
\usepackage{cite}
\usepackage{stfloats}
\usepackage{url}
\usepackage{listings}
\usepackage{fancyhdr}
\usepackage{biblatex}
\usepackage{color}
\addbibresource{sources.bib}
%\usepackage{hyperref}
%\usepackage{wrapfig}
\usepackage{array}
%\usepackage{multirow}
%\usepackage{adjustbox}
%\usepackage{nccmath}
\usepackage[dvipsnames]{xcolor}
\usepackage{tcolorbox}
\newtcolorbox{mybox}{colback=white,
colframe=black}
\newcommand{\titulo}{Abgabe 3; Diskrete Mathematik}
\newcommand{\nombre}{Liliana Sanfilippo}
\newcommand{\FCE}{\textbf{WiSe 2024/2025}}

 

\usepackage{fancyhdr}
\pagestyle{fancy}
\fancyhead{}
\fancyfoot{}
\fancyhead[LO]{\small\nombre}
\fancyhead[RE]{\small\titulo}
\fancyhead[LO]{\small\FCE}
\fancyfoot[RO,LE] {\thepage}

\setcounter{page}{1} %Número de la página, la editorial lo cambiará



\usepackage[onehalfspacing]{setspace}
\begin{document}


\begin{tabular}{p{12cm}}
{{\Large\bf\titulo}}\\
\vspace{0.2cm}\\
 {\large\nombre}\\
 \vspace{0.2cm}\\   
\end{tabular}
\section{Präsenzaufgaben, im Tutorium bearbeitet}
\subsection{Aufgabe 1}
Die Permutationen $\pi$ und $\sigma$ von $\{1, 2, \dots, 7\}$ seien in Zykelschreibweise durch
    \[
    \pi := (1, 2)(3, 4, 5, 6)(7) \quad \text{und} \quad \sigma := (1)(2, 3, 4, 5, 6, 7) 
    \]
    
    gegeben. Stellen Sie $\pi\sigma$, $\sigma\pi$, $\pi^2$, $\pi^{-1}$ und $\sigma^{-1}$ als Produkte disjunkter Zykeln dar.
    \subsection*{$\pi\sigma$:}
    \begin{align*}
        \pi (\sigma (1)) = \pi (1) &  =2 & \text{also } 1 \xrightarrow[]{\pi \sigma} 2 \\
        \pi (\sigma (2)) = \pi (3) &  =4 & \text{also } 2 \xrightarrow[]{\pi \sigma} 4 \\
        \pi (\sigma (3)) = \pi (4) &  =5 & \text{also } 3 \xrightarrow[]{\pi \sigma} 5 \\
        \pi (\sigma (4)) = \pi (5) &  = 6& \text{also } 4 \xrightarrow[]{\pi \sigma} 6 \\
        \pi (\sigma (5)) = \pi (6) &  =3 & \text{also } 5 \xrightarrow[]{\pi \sigma} 3 \\
        \pi (\sigma (6)) = \pi (7) &  =7 & \text{also } 6 \xrightarrow[]{\pi \sigma} 7 \\
        \pi (\sigma (7)) = \pi (2) &  = 1& \text{also } 7 \xrightarrow[]{\pi \sigma} 1 \\
    \end{align*}
Als Produkt disjunkter Zykel: 
\[
(3,5)(1,2,4,6,7)\textcolor{Blue}{^*}
\]
\textcolor{Blue}{$^*$ Bei disjunkten Zykeln: (3,5)(1,2,4,6,7) = (1,2,4,6,7)(3,5)} \newline
\subsection*{$\sigma \pi$:}
 \begin{align*}
        \sigma (\pi (1)) = \sigma (2) &  =3 & \text{also } 1 \xrightarrow[]{\sigma \pi} 3 \\
        \sigma (\pi (2)) = \sigma (1) &  =1 & \text{also } 2 \xrightarrow[]{\sigma \pi} 1 \\
        \sigma (\pi (3)) = \sigma (4) &  =5 & \text{also } 3 \xrightarrow[]{\sigma \pi} 5 \\
        \sigma (\pi (4)) = \sigma (5) &  =6 & \text{also } 4 \xrightarrow[]{\sigma \pi} 6 \\
        \sigma (\pi (5)) = \sigma (6) &  =7 & \text{also } 5 \xrightarrow[]{\sigma \pi} 7 \\
        \sigma (\pi (6)) = \sigma (3) &  =4 & \text{also } 6 \xrightarrow[]{\sigma \pi} 4 \\
        \sigma (\pi (7)) = \sigma (7) &  =2 & \text{also } 7 \xrightarrow[]{\sigma \pi} 2 \\
    \end{align*}
Als Produkt disjunkter Zykel: 
\[
(1,3,5,7,2)(4,6)
\]
\subsection*{$\pi^2$:}
\begin{align*}
     \sigma (\sigma (1)) = \sigma (2) &  =1 & \text{also } 1 \xrightarrow[]{\sigma \sigma} 1 \\
     \sigma (\sigma (2)) = \sigma (1) &  =2 & \text{also } 2 \xrightarrow[]{\sigma \sigma} 2 \\
     \sigma (\sigma (3)) = \sigma (4) &  =5 & \text{also } 3 \xrightarrow[]{\sigma \sigma} 5 \\
     \sigma (\sigma (4)) = \sigma (5) &  =6 & \text{also } 4 \xrightarrow[]{\sigma \sigma} 6 \\
     \sigma (\sigma (5)) = \sigma (6) &  =3 & \text{also } 5 \xrightarrow[]{\sigma \sigma} 3 \\
     \sigma (\sigma (6)) = \sigma (3) &  =4 & \text{also } 6 \xrightarrow[]{\sigma \sigma} 4   \\
     \sigma (\sigma (7)) = \sigma (7) &  = 7 & \text{also } 7 \xrightarrow[]{\sigma \sigma} 7 \\
\end{align*}
Als Produkt disjunkter Zykel: 
\[
(3,5)(4,6)
\]
\subsection*{$\pi^{-1}$:}
 \[
    \pi^{-1} := (1, 2)(6,5,4,3) 
    \]
 \subsection*{$\sigma^{-1}$:}   
    \[
    \sigma^{-1} := (7,6,5,4,3,2)
    \]
\textcolor{Blue}{\textit{Anmerkung: \\ Die Fixpunkte weglassen, weil das keine Zykel sind.}}
\subsection{Aufgabe 2}
     Seien $A$, $B$, $C$, $D \subseteq X$. Geben Sie die Formel für $|A \cup B \cup C \cup D|$ in Abhängigkeit der Mächtigkeit der Durchschnitte an.

\[
|A \cup B \cup C \cup D| = |A| + |B| + |C| + |D| 
\]
\[
- |A \cap B| - |A \cap C| - |A \cap D| - |B \cap C| - |B \cap D| - |C \cap D|
\]
\[
+ |A \cap B \cap C| + |A \cap B \cap D| + |A \cap C \cap D| + |B \cap C \cap D|
\]
\[
- |A \cap B \cap C \cap D|
\]
\color{Blue}
\subsubsection*{\textit{Notizen für eigenes Verständnis:}}
Das ist die Inklusions-Exklusions-Formel nur für vier Mengen. Sie ergibt die Anzahl der Elemente der Vereinigung und sorgt dafür, dass Sachen aber nicht doppelt gezählt werden. 
\begin{table}[h]
\color{Blue}
    \centering
    \begin{tabular}{p{8cm} p{9cm}}
        $|A| + |B| + |C| + |D| $ & alle Elemente der Mengen zusammen addieren \\ \\
        $- |A \cap B| - |A \cap C| - |A \cap D| - |B \cap C| - |B \cap D| - |C \cap D|$ &  Paarweise Überschneidungen entfernen \\ \\
        $+ |A \cap B \cap C| + |A \cap B \cap D| + |A \cap C \cap D| + |B \cap C \cap D|$ & Dreifache Überschneidungen wieder hinzufügen, weil man mit den Paaren zu viel weg gemacht hat \\ \\
        $|A \cap B \cap C \cap D|$ & Vierfache Überschneidung wegmachen \\
    \end{tabular}
\end{table}
\color{black}
\section{Aufgaben zur Abgabe}
\subsection*{Aufgabe 1 (4 Punkte)}
Eine Permutation $\pi$ mit der Eigenschaft $\pi(i) \neq i$ für alle $i \in \{1, 2, \dots, n\}$ heißt
\textit{Derangement} (oder auch fixpunktfreie Permutation \textcolor{Blue}{d.h. Permutationen, bei denen kein Element auf seiner ursprünglichen Position bleibt}). Wir bezeichnen die Teilmenge der
Derangements mit $D_n \subseteq S_n$.

Sei $X = S_n$ und $A_i = \{\pi \in S_n \mid \pi(i) = i\}$ für $1 \leq i \leq n$. Für $I \subseteq \{1, 2, \dots, n\}$ schreiben wir
$A_I = \bigcap_{i \in I} A_i$ \textcolor{Blue}{(d.h. $A_i$ ist die Menge aller Permutationen, bei denen das i-te Element an seiner ursprünglichen Position bleibt und $A_I$ besteht aus den Permutationen $\pi$, bei denen alle Elemente an den Positionen $i \in I$ Fixpunkte sind - also nicht nur mindestens eines!)} \\
\textcolor{Blue}{\textbf{Beispiel: }Wenn \( I = \{1, 2\} \), dann ist \( A_I = A_1 \cap A_2 \), und es umfasst alle Permutationen, bei denen sowohl \( \pi(1) = 1 \) als auch \( \pi(2) = 2 \).
}
\begin{enumerate}
    \item[(a)] Finden Sie $|A_I|$ für $I \subseteq \{1, 2, \dots, n\}$. 
   % \begin{mybox}
    %    \[
     %   |A_I| = \sum_{k=1}^{|I|} (-1)^{k+1} \binom{|I|}{k} (n - k)! 
     %   \] 
      %  Dabei ist $k = |J|$, also die Größe einer Menge $A_J \subseteq A_I$. \\
       %  Man möchte Mehrfachzählungen vermeiden und daher Subtrahiert und Addiert man abwechselnd die z-fachen Überlappungen, deren Anzahl jeweils der der k-elementigen Teilmengen $J \subseteq I$ entspricht. Dabei ist $z \in \{1, ..., |I|, |I|+1\}$, also $z$ rangiert von $1$ bis $|I|+1$. Deswegen steht $(-1)^{k+1}$ in der Formel. \\
     %    Da wir $|A_I|$ berechnen wollen, zählen wir natürlich wie viele z-fache Überlappungen bzw. k-elementigen Teilmengen es pro Iteration in I gibt. Aus Definition 2.1. wissen wir, es gibt $\binom{|I|}{k}$ k-elementige Teilmengen in einer $|I|$-elementigen Menge. \\
     %    $(n-k)!$ bezeichnet die Anzahl der nicht fixierten Punkte, welche ihre Positionen verändern könnten und somit die Anzahl der möglichen Permutationen erhöhen.  
    %    \textcolor{white}{\cite{inkludionsprinzipENG}}
%\end{mybox}
\begin{mybox}
    Wenn alle Elemente in $I$ fixiert sind, können immer noch $n - |I|$ Elemente frei permuttiert werden. Aus der Vorlesung wissen wir, dass $(n - |I|)!$ mögliche Permutationen ergibt. \\
    Daher: 
    \[
        A_I = (n - |I|)!
    \]
\end{mybox}
    \item[(b)] Finden Sie $|D_n|$.
\begin{mybox}
    Bei $D_n$ handelt es sich um eine Teilmenge von $S_n$, in der keine Permutation Fixpunkte aus $I$ hat. Also auch nicht mindestens einen.  \\
    Die Permutationen mit mindestens einem fixierten Punkt können durch eine Teilmenge 
    \[
       |B_I| = \sum_{k=1}^{|I|} (-1)^{k+1} \binom{|I|}{k} (n - k)! 
    \] 
    dargestellt werden. 
    Dabei ist $k = |J|$, also die Größe einer Menge $A_J \subseteq A_I$. \\
    Man möchte Mehrfachzählungen vermeiden und daher Subtrahiert und Addiert man abwechselnd die z-fachen Überlappungen, deren Anzahl jeweils der der k-elementigen Teilmengen $J \subseteq I$ entspricht. Dabei ist $z \in \{1, ..., |I|, |I|+1\}$, also $z$ rangiert von $1$ bis $|I|+1$. Deswegen steht $(-1)^{k+1}$ in der Formel. \\
    Da wir $|A_I|$ berechnen wollen, zählen wir natürlich wie viele z-fache Überlappungen bzw. k-elementigen Teilmengen es pro Iteration in I gibt. Aus Definition 2.1. wissen wir, es gibt $\binom{|I|}{k}$ k-elementige Teilmengen in einer $|I|$-elementigen Menge. \\
    $(n-k)!$ bezeichnet die Anzahl der nicht fixierten Punkte, welche ihre Positionen verändern könnten und somit die Anzahl der möglichen Permutationen erhöhen. \\
    $D_n$ ist das Komplement von $B_I$.   
    \[
    \text{Also: } D_n = S_n \backslash B_I
    \]
    \begin{align*}
         \text{Und: } |D_n| & = |S_n| - |B_I| \\
        & = |S_n| - \sum_{k=1}^{|I|} (-1)^{k+1} \binom{|I|}{k} (n - k)! 
    \end{align*}
    
\end{mybox}
\end{enumerate}


\subsection{Aufgabe 2 (4 Punkte)}
In einer Gruppe von 70 Studierenden sprechen 45 Englisch, 37 Französisch und 22 beide
Sprachen. Wie viele sprechen keine von beiden Sprachen?
\begin{mybox}
\begin{align*}
    45 + 37 - 22 & = 60 \\
    72 - 60 & = 12
\end{align*}
12 Studierende sprechen keine von beiden Sprachen. 
\end{mybox}
\noindent Wenn außerdem 16 Studierende Russisch sprechen, von denen 5 auch Englisch, 7 auch
Französisch und 4 alle drei Sprachen sprechen, wie viele sprechen dann keine der drei
Sprachen?
\begin{mybox}
\begin{align*}
    45 + 37 + 16 & = 98 \\
    98 - 22 - 5 - 7 & = 64 \\
    64 + 4 & = 68 \\ 
    70 - 68 & = 2 
\end{align*}
2 Studierende sprechen keine der drei Sprachen.
\end{mybox}

\subsection*{Aufgabe 3 (4 Punkte)}
Ein Kartenspiel mit $2n$ Karten wird wie folgt gemischt:
Der Stapel wird in obere und untere Hälfte geteilt. Die beiden Hälften werden dann so
ineinander geschoben, dass der neu entstehende Stapel von oben beginnend abwechselnd
eine Karte von der unteren und der oberen Hälfte enthält. \\
Für die gegebenen $n$, wie oft muss gemischt werden, bis sich die Karten wieder in ihrer
Ursprungsreihenfolge befinden? \\
\textcolor{Blue}{Der Mischvorgang ist eine Permutation. Da nach jeder Permutation jeweils eine Karte aus dem oberen Stapel sich mit einer aus dem unteren Stapel abwechselt, ist die Permutation der Art $\pi = (1,n+1)(2,n+2)...$}
\begin{enumerate}
    \item[(a)] $n = 2$
    \begin{mybox}
    $ 2 \cdot 2 = 4$ \\ \newline
        \begin{tabular}{c|cccc}
            Mischvorgang 0 & 1 & 2& 3& 4\\
            Mischvorgang 1 & 1 & 3 & 2 & 4 \\
            Mischvorgang 2 & 1 & 2 & 3 & 4
        \end{tabular} \\ \newline
    Es braucht 2 Mischungen, damit die Karten wieder in der richtigen Reihenfolge sind. 
    \end{mybox}
    \item[(b)] $n = 4$
    \begin{mybox}
     $ 2 \cdot 4 = 8$ \\ \newline
     \begin{tabular}{c|cccccccc}
        Mischvorgang 0  &1&2&3&4&5&6&7&8 \\
        Mischvorgang 1  &1&5&2&6&3&7&4&8  \\
        Mischvorgang 2 &1&3&5&7&2&4&6&8  \\
        Mischvorgang 3 &1&2&3&4&5&6&7&8  \\
     \end{tabular} \\ \newline
     Es braucht 3 Mischungen, damit die Karten wieder in der richtigen Reihenfolge sind.
    \end{mybox}
    \item[(c)] $n = 8$
\begin{mybox}
  $ 2 \cdot 8 = 16$ \\ \newline
     \begin{tabular}{c|cccccccccccccccc}
        Mischvorgang 0  &1&2&3&4&5&6&7&8&9&10&11&12&13&14&15&16 \\
        Mischvorgang 1  &1&9&2&10&3&11&4&12&5&13&6&14&7&15&8&16  \\
        Mischvorgang 2 &1&5&9&13&2&6&10&14&3&7&11&15&4&8&12&16  \\
        Mischvorgang 3 &1&3&5&7&9&11&13&15&2&4&6&8&10&12&14&16  \\
        Mischvorgang 4 &1&2&3&4&5&6 &7 & 8&9&10&11&12&13&14&15&16  \\
     \end{tabular} \\ \newline
     Es braucht 4 Mischungen, damit die Karten wieder in der richtigen Reihenfolge sind.
\end{mybox}
\end{enumerate}


\subsection*{Aufgabe 4 (4 Punkte)} 
Gegeben sei eine Permutation $\sigma$ der Menge $\{1, 2, \dots, n\}$. \\
Zeigen Sie, dass $\operatorname{sgn}(\sigma) = \det(M(\sigma))$, wobei $M(\sigma)$ die $(n \times n)$-Matrix bezeichnet, deren $i$-te
Spalte der $\sigma(i)$-te Standardbasisvektor des $\mathbb{R}^n$ ist.
\begin{mybox}
  Aufgrund der gegebenen Definition von $M(\sigma)$, kann man es so ausschreiben: $M(\sigma)_{i,j} = 
\begin{cases} 
1 & \text{wenn } j = \sigma(i), \\
0 & \text{sonst}.
\end{cases}$ \\
Dank Leibniz wissen wir, dass 
\[
\det(A) = \sum_{\sigma \in S_n} \text{sgn}(\sigma) \prod_{i=1}^n a_{i, \sigma(i)}
\] 
also hier \[
\det(M(\sigma)) = \sum_{\sigma \in S_n} \text{sgn}(\sigma) \prod_{i=1}^n M(\sigma)_{i, \sigma(i)}
\]
und weil \( M(\sigma)_{i, \sigma(i)} = 1 \) für jedes \( i \) ist es \[
\det(M(\sigma)) = \sum_{\sigma \in S_n} \text{sgn}(\sigma) \cdot 1
\]
Also: 
\[
\det(M(\sigma)) =  \text{sgn}(\sigma) 
\]

\end{mybox}


\end{document}